\documentclass[11pt]{letter}
\usepackage[utf8]{inputenc}
\usepackage[T1]{fontenc}
\usepackage[margin=1in]{geometry}
\usepackage{hyperref}

\signature{Richar Andre Vilca-Solorzano\\
  on behalf of all co-authors\\
  Universidad Nacional del Altiplano (UNAP), Puno, Peru\\
  \href{mailto:75521963@est.unap.edu.pe}{75521963@est.unap.edu.pe}}

\address{Faculty of Statistical and Computer Engineering\\
  Universidad Nacional del Altiplano (UNAP)\\
  Puno, Peru}

\begin{document}

\begin{letter}{The Editor-in-Chief\\ \textit{Natural Hazards}\\ Springer}

\opening{Dear Editor,}

We are pleased to submit our original research manuscript entitled
\textbf{``Pan-Andean explainable machine learning identifies priority
glacial lakes for outburst flood early warning''} for consideration for
publication in \textit{Natural Hazards}.

Glacial lake outburst floods (GLOFs) are among the most destructive
cryospheric hazards in the Andes, yet disaster-risk managers lack an
objective, scalable method to decide which of the tens of thousands of
Andean glacial lakes most urgently require inspection. Our study
addresses this gap directly and, to our knowledge, presents the
\emph{first pan-Andean, explainable GLOF susceptibility model} validated
against an independent historical catalogue. The main contributions are:

\begin{itemize}
\item A harmonized, multitemporal glacial-lake inventory of 69{,}643
  lake-year observations (20{,}140 distinct lakes $\geq$0.05 km$^2$)
  spanning fifteen study areas across four countries (Peru, Bolivia,
  Ecuador, and Chile, including Patagonia), built from free Sentinel-2
  and Landsat imagery and open digital elevation models.
\item A rigorously validated, \emph{leakage-free} machine-learning
  framework evaluated under strict StratifiedGroup out-of-fold and
  Leave-One-Country-Out cross-validation, explicitly designed for the
  extreme class imbalance (1:629) and data scarcity ($n=34$ GLOF
  lake-years from 32 distinct source lakes) inherent to historical GLOF
  catalogues. The best discriminator reaches AUC-ROC 0.826
  (95\% CI 0.729--0.916; permutation $p<0.001$).
\item SHAP-based physical interpretation that yields an honest,
  multivariate signature of GLOF-prone lakes---moderate size combined
  with compact shape, low moraine-dam freeboard, and proximity to active
  ice---rather than a single-variable proxy.
\item A fully open-source, reproducible pipeline and an operational,
  isotonically calibrated (Brier Skill Score 0.24) prioritization
  framework: a compact top-2\% watch-list of 403 lakes recovers 41\% of
  confirmed GLOF-source lakes at $20.3\times$ random enrichment
  ($74.8\times$ in the top 0.5\%), directly usable by resource-constrained
  agencies such as Peru's INAIGEM for evidence-based disaster risk
  reduction.
\end{itemize}

We believe the manuscript is well aligned with the scope of
\textit{Natural Hazards}, which covers flood and glacial hazards,
risk assessment, and disaster risk reduction. Critically, the framework
is designed for impact in \emph{developing Andean nations} that bear the
highest GLOF exposure yet have the fewest resources: it relies only on
free Sentinel-2 imagery and open DEMs, runs without proprietary
software, and turns these inputs into a ranked watch-list that lets
resource-constrained agencies (e.g.\ Peru's INAIGEM) triage scarce
field-inspection budgets toward the lakes most likely to threaten
downstream communities. The work should thus be of broad interest to the
hazard, cryosphere, remote-sensing, and disaster-risk-reduction
communities, and provides a transferable template for other glacierized
mountain regions.

We confirm that this manuscript is original, has not been published
previously, and is not under consideration for publication elsewhere.
All authors have read and approved the submission and agree to its
publication. The authors declare no competing interests. All input
datasets are publicly available, and the full analysis pipeline,
trained models, and code are openly released under the MIT licence.

We suggest the \emph{Natural Hazards and Risk} subject area and would be
glad to recommend potential reviewers upon request. Thank you for your
consideration.

\closing{Sincerely,}

\end{letter}
\end{document}
